\documentclass[border=4pt]{standalone}
\usepackage{amsmath,amssymb,mathtools,xcolor,varwidth}
\begin{document}
\begin{varwidth}{28em}
\large\bfseries
The alteration is made in the direction of the right line in which that force is impress'd — always directed the same way with the generating force. If the body moved before, this new motion is added to or subducted from the former motion, according as they directly conspire or are directly contrary. When they are obliquely joyned, the force and the prior motion compound by the parallelogram rule, so as to produce a new motion compounded from the determination of both — hence $\Delta(\text{motion})\ \parallel F\$.
\end{varwidth}
\end{document}
